
\documentclass{article}
\usepackage{colortbl}
\usepackage{makecell}
\usepackage{multirow}
\usepackage{supertabular}

\begin{document}

\newcounter{utterance}

\centering \large Interaction Transcript for game `dond', experiment `semi', episode 2 with gpt{-}5.4{-}azure.
\vspace{24pt}

{ \footnotesize  \setcounter{utterance}{1}
\setlength{\tabcolsep}{0pt}
\begin{supertabular}{c@{$\;$}|p{.15\linewidth}@{}p{.15\linewidth}p{.15\linewidth}p{.15\linewidth}p{.15\linewidth}p{.15\linewidth}}
    \# & $\;$A & \multicolumn{4}{c}{Game Master} & $\;\:$B\\
    \hline

    \theutterance \stepcounter{utterance}  
    & & \multicolumn{4}{p{0.6\linewidth}}{
        \cellcolor[rgb]{0.9,0.9,0.9}{
            \makecell[{{p{\linewidth}}}]{
                \texttt{\tiny{[P1$\langle$GM]}}
                \texttt{Hráš kooperatívnu hru rokovania, v ktorej sa musíš dohodnúť s iným hráčom na tom, ako si rozdeliť súbor položiek medzi sebou.} \\
\\ 
\texttt{Pravidlá:} \\
\texttt{(a) Ty a druhý hráč dostávate sadu položiek. Každý z vás dostane aj tajnú funkciu hodnoty, ktorá predstavuje, ako si ceníš jednotlivé typy objektov.} \\
\texttt{(b) Vymieňate si správy s druhým hráčom, aby ste sa dohodli na tom, kto dostane ktoré položky. Každý môže poslať maximálne 5 správ alebo ukončiť hru skôr tým, že kedykoľvek predložíte tajný návrh.} \\
\texttt{(c) Od vás sa vyžaduje predložiť tajný návrh označujúci položky, ktoré chcete, formátovaný v hranatých zátvorkách takto: "{[}Návrh: <číslo> <názov objektu>, <číslo> <názov objektu>, <...>{]}"} \\
\texttt{(d) Ak sú vaše návrhy komplementárne, t.j. je dostatok položiek na splnenie oboch návrhov, každý hráč dostane skóre na základe súčtu hodnôt položiek, ktoré dostal. V opačnom prípade dostávate obaja nula bodov.} \\
\texttt{(e) Vaším cieľom je maximalizovať súčet vášho skóre a skóre druhého hráča.} \\
\\ 
\texttt{Začnime.} \\
\\ 
\texttt{Sada dostupných položiek je:} \\
\\ 
\texttt{1 taška, 3 lampy, 1 kvet} \\
\\ 
\texttt{Vaša tajná funkcia hodnoty je:} \\
\\ 
\texttt{taška: 3, lampa: 0, kvet: 7} \\
\\ 
\texttt{DÔLEŽITÉ: Vaše správy, pokiaľ nejde o tajný návrh, sa bezprostredne prepošlú druhému hráčovi, preto nezahrňujte žiadnu odpoveď na pravidlá ani text oznamujúci vašu správu. Na podanie tajného návrhu použite uvedený formát. Nepoužívajte hranaté zátvorky, keď komunikujete s druhým hráčom, inak sa to bude interpretovať ako váš tajný návrh.} \\
            }
        }
    }
    & \\ \\

    \theutterance \stepcounter{utterance}  
    & \multicolumn{4}{p{0.6\linewidth}}{
        \cellcolor[rgb]{0.9,0.9,0.9}{
            \makecell[{{p{\linewidth}}}]{
                \texttt{\tiny{[P1$\rangle$GM]}}
                \texttt{Mne na lampách nezáleží, najviac chcem kvet, taška je pre mňa trochu hodnotná. Ak chceš, môžeš si vziať všetky 3 lampy a potom doladíme tašku a kvet podľa toho, čo je cennejšie pre teba.} \\
            }
        }
    }
    & & \\ \\

    \theutterance \stepcounter{utterance}  
    & & & \multicolumn{2}{p{0.3\linewidth}}{
        \cellcolor[rgb]{0.9,0.9,0.9}{
            \makecell[{{p{\linewidth}}}]{
                \texttt{\tiny{[GM$|$GM]}}
                \texttt{continue} \\
            }
        }
    }
    & & \\ \\

    \theutterance \stepcounter{utterance}  
    & & & \multicolumn{2}{p{0.3\linewidth}}{
        \cellcolor[rgb]{0.9,0.9,0.9}{
            \makecell[{{p{\linewidth}}}]{
                \texttt{\tiny{[GM$|$GM]}}
                \texttt{Mne na lampách nezáleží, najviac chcem kvet, taška je pre mňa trochu hodnotná. Ak chceš, môžeš si vziať všetky 3 lampy a potom doladíme tašku a kvet podľa toho, čo je cennejšie pre teba.} \\
            }
        }
    }
    & & \\ \\

    \theutterance \stepcounter{utterance}  
    & & \multicolumn{4}{p{0.6\linewidth}}{
        \cellcolor[rgb]{0.9,0.9,0.9}{
            \makecell[{{p{\linewidth}}}]{
                \texttt{\tiny{[GM$\rangle$P2]}}
                \texttt{Hráš kooperatívnu hru rokovania, v ktorej sa musíš dohodnúť s iným hráčom na tom, ako si rozdeliť súbor položiek medzi sebou.} \\
\\ 
\texttt{Pravidlá:} \\
\texttt{(a) Ty a druhý hráč dostávate sadu položiek. Každý z vás dostane aj tajnú funkciu hodnoty, ktorá predstavuje, ako si ceníš jednotlivé typy objektov.} \\
\texttt{(b) Vymieňate si správy s druhým hráčom, aby ste sa dohodli na tom, kto dostane ktoré položky. Každý môže poslať maximálne 5 správ alebo ukončiť hru skôr tým, že kedykoľvek predložíte tajný návrh.} \\
\texttt{(c) Od vás sa vyžaduje predložiť tajný návrh označujúci položky, ktoré chcete, formátovaný v hranatých zátvorkách takto: "{[}Návrh: <číslo> <názov objektu>, <číslo> <názov objektu>, <...>{]}"} \\
\texttt{(d) Ak sú vaše návrhy komplementárne, t.j. je dostatok položiek na splnenie oboch návrhov, každý hráč dostane skóre na základe súčtu hodnôt položiek, ktoré dostal. V opačnom prípade dostávate obaja nula bodov.} \\
\texttt{(e) Vaším cieľom je maximalizovať súčet vášho skóre a skóre druhého hráča.} \\
\\ 
\texttt{Začnime.} \\
\\ 
\texttt{Sada dostupných položiek je:} \\
\\ 
\texttt{1 taška, 3 lampy, 1 kvet} \\
\\ 
\texttt{Vaša tajná funkcia hodnoty je:} \\
\\ 
\texttt{taška: 1, lampa: 2, kvet: 3} \\
\\ 
\texttt{DÔLEŽITÉ: Vaše správy, pokiaľ nejde o tajný návrh, sa bezprostredne prepošlú druhému hráčovi, preto nezahrňujte žiadnu odpoveď na pravidlá ani text oznamujúci vašu správu. Na podanie tajného návrhu použite uvedený formát. Nepoužívajte hranaté zátvorky, keď komunikujete s druhým hráčom, inak sa to bude interpretovať ako váš tajný návrh.} \\
\\ 
\texttt{Mne na lampách nezáleží, najviac chcem kvet, taška je pre mňa trochu hodnotná. Ak chceš, môžeš si vziať všetky 3 lampy a potom doladíme tašku a kvet podľa toho, čo je cennejšie pre teba.} \\
            }
        }
    }
    & \\ \\

    \theutterance \stepcounter{utterance}  
    & & & \multicolumn{4}{p{0.6\linewidth}}{
        \cellcolor[rgb]{0.9,0.9,0.9}{
            \makecell[{{p{\linewidth}}}]{
                \texttt{\tiny{[GM$\langle$P2]}}
                \texttt{To mi sedí — vezmem si všetky 3 lampy a ty si nechaj kvet aj tašku.} \\
            }
        }
    }
     \\ \\

    \theutterance \stepcounter{utterance}  
    & & & \multicolumn{2}{p{0.3\linewidth}}{
        \cellcolor[rgb]{0.9,0.9,0.9}{
            \makecell[{{p{\linewidth}}}]{
                \texttt{\tiny{[GM$|$GM]}}
                \texttt{continue} \\
            }
        }
    }
    & & \\ \\

    \theutterance \stepcounter{utterance}  
    & & & \multicolumn{2}{p{0.3\linewidth}}{
        \cellcolor[rgb]{0.9,0.9,0.9}{
            \makecell[{{p{\linewidth}}}]{
                \texttt{\tiny{[GM$|$GM]}}
                \texttt{To mi sedí — vezmem si všetky 3 lampy a ty si nechaj kvet aj tašku.} \\
            }
        }
    }
    & & \\ \\

    \theutterance \stepcounter{utterance}  
    & & \multicolumn{4}{p{0.6\linewidth}}{
        \cellcolor[rgb]{0.9,0.9,0.9}{
            \makecell[{{p{\linewidth}}}]{
                \texttt{\tiny{[P1$\langle$GM]}}
                \texttt{To mi sedí — vezmem si všetky 3 lampy a ty si nechaj kvet aj tašku.} \\
            }
        }
    }
    & \\ \\

    \theutterance \stepcounter{utterance}  
    & \multicolumn{4}{p{0.6\linewidth}}{
        \cellcolor[rgb]{0.9,0.9,0.9}{
            \makecell[{{p{\linewidth}}}]{
                \texttt{\tiny{[P1$\rangle$GM]}}
                \texttt{{[}Návrh: 1 taška, 1 kvet{]}} \\
            }
        }
    }
    & & \\ \\

    \theutterance \stepcounter{utterance}  
    & & & \multicolumn{2}{p{0.3\linewidth}}{
        \cellcolor[rgb]{0.9,0.9,0.9}{
            \makecell[{{p{\linewidth}}}]{
                \texttt{\tiny{[GM$|$GM]}}
                \texttt{proposal} \\
            }
        }
    }
    & & \\ \\

    \theutterance \stepcounter{utterance}  
    & & & \multicolumn{2}{p{0.3\linewidth}}{
        \cellcolor[rgb]{0.9,0.9,0.9}{
            \makecell[{{p{\linewidth}}}]{
                \texttt{\tiny{[GM$|$GM]}}
                [1, 0, 1]
            }
        }
    }
    & & \\ \\

    \theutterance \stepcounter{utterance}  
    & & \multicolumn{4}{p{0.6\linewidth}}{
        \cellcolor[rgb]{0.9,0.9,0.9}{
            \makecell[{{p{\linewidth}}}]{
                \texttt{\tiny{[GM$\rangle$P2]}}
                \texttt{Druhý hráč predložil tajný návrh. Teraz je na vás, aby ste predložili vlastný návrh.} \\
            }
        }
    }
    & \\ \\

    \theutterance \stepcounter{utterance}  
    & & & \multicolumn{4}{p{0.6\linewidth}}{
        \cellcolor[rgb]{0.9,0.9,0.9}{
            \makecell[{{p{\linewidth}}}]{
                \texttt{\tiny{[GM$\langle$P2]}}
                \texttt{{[}Návrh: 0 taška, 3 lampy, 0 kvet{]}} \\
            }
        }
    }
     \\ \\

    \theutterance \stepcounter{utterance}  
    & & & \multicolumn{2}{p{0.3\linewidth}}{
        \cellcolor[rgb]{0.9,0.9,0.9}{
            \makecell[{{p{\linewidth}}}]{
                \texttt{\tiny{[GM$|$GM]}}
                \texttt{proposal} \\
            }
        }
    }
    & & \\ \\

    \theutterance \stepcounter{utterance}  
    & & & \multicolumn{2}{p{0.3\linewidth}}{
        \cellcolor[rgb]{0.9,0.9,0.9}{
            \makecell[{{p{\linewidth}}}]{
                \texttt{\tiny{[GM$|$GM]}}
                [0, 3, 0]
            }
        }
    }
    & & \\ \\

    \theutterance \stepcounter{utterance}  
    & & & \multicolumn{2}{p{0.3\linewidth}}{
        \cellcolor[rgb]{0.9,0.9,0.9}{
            \makecell[{{p{\linewidth}}}]{
                \texttt{\tiny{[GM$|$GM]}}
                [[1, 0, 1], [0, 3, 0]]
            }
        }
    }
    & & \\ \\

\end{supertabular}
}

\end{document}
